\section{Concluding Remarks}

This proposal has set out a retrieval-augmented, semi-parametric continual-learning framework for Vietnamese legal language models. Its necessity arises from a dual challenge---keeping legal knowledge current while forgetting specific information on demand---that conventional weight-updating approaches cannot meet without catastrophic forgetting or costly, hard-to-verify unlearning. The challenge is timely: legislation and case law change continuously, data-governance obligations such as the right to be forgotten are tightening, and high-quality Vietnamese legal benchmarks have only recently become available, together providing both the motivation and the empirical foundation for the study.

The framework is also feasible for the author to undertake. Prior hands-on experience designing a retrieval-augmented generation system during an undergraduate thesis (on university-information question answering) provides the technical foundation, which this work transfers and extends to the substantially more demanding legal domain---with its temporal validity, citation structure, and compliance constraints---under the supervision of an advisor specializing in natural language processing and drawing on publicly available Vietnamese legal datasets.

By freezing the base model and confining all knowledge updates and unlearning to an external knowledge base, the proposed approach delivers zero-weight-drift knowledge maintenance, compliance-aware controllable forgetting, and a reproducible evaluation protocol. These contributions are intended to be both scientifically novel---bridging continual learning, retrieval-augmented generation, and legal temporal regimes for a low-resource language---and practically valuable for Vietnam's emerging legal-artificial-intelligence ecosystem.
